\subsection{Bregman 近端中的核平坦性与物理速率}
\label{subsec:gppa-bregman-quartic}

Bregman 近端法为非线性核提供了一个直接的优化解释。
对全空间上的可微凸函数 $\phi$，记
\[
 D_\phi(y,x)=\phi(y)-\phi(x)-\langle\nabla\phi(x),y-x\rangle.
\]
对于适当下半连续凸函数 $f$，最小化 $\lambda f(y)+D_\phi(y,x)$ 的最优性条件为
$\nabla\phi(x)\in\nabla\phi(y)+\lambda\partial f(y)$
\cite{BauschkeBorweinCombettes2003}。
下面的二次目标已足以说明，核在解附近的平坦性能够改变物理迭代的速率。

在 $\mathbb R^2$ 上取
\[
 f(t,s)=\frac12t^2,\qquad
 \phi(t,s)=\frac14t^4+\frac12s^2,
 \qquad v(t,s)=\nabla\phi(t,s)=(t^3,s).
\]
目标的解集是直线 $S=\{0\}\times\mathbb R$。
核 $v$ 是全空间上的非线性同胚，但其逆在 $S$ 附近不局部 Lipschitz。
完整 Bregman 近端步满足
\begin{equation}
\label{eq:gppa-bregman-cubic-step}
 t_k^3=t_{k+1}^3+\lambda t_{k+1},\qquad
 s_{k+1}=s_k.
\end{equation}
因为 $u\mapsto u^3+\lambda u$ 连续、严格递增且满射，
这一完整近端映射对每个输入都存在且唯一。

\begin{proposition}
\label{prop:gppa-bregman-cubic}
对任意 $\lambda>0$ 与任意初值 $(t_0,s_0)$，
迭代 \eqref{eq:gppa-bregman-cubic-step} 收敛到 $(0,s_0)$，
物理轨道的总长度恰为 $|t_0|$。
若 $t_0\ne0$，令 $e_k=|t_k|$，则
\[
 \lim_{k\to\infty}\frac{e_{k+1}}{e_k^3}=\frac1\lambda.
\]
\end{proposition}
\begin{proof}
若 $t_0=0$，轨道驻定。若 $t_0\ne0$，
方程~\eqref{eq:gppa-bregman-cubic-step} 使 $t_{k+1}$ 与 $t_k$ 同号，
且 $|t_{k+1}|<|t_k|$，故 $(t_k)$ 有极限。
在方程中取极限可知 $t_k\to0$，而 $s_k=s_0$。
于是总长度由 $\sum_k|t_k-t_{k+1}|=|t_0|$ 给出。
同一方程还给
\[
 \frac{|t_{k+1}|}{|t_k|^3}
 =\frac1{\lambda+t_{k+1}^2}\longrightarrow\frac1\lambda.
\]
\end{proof}

本例对固定核定理~\ref{thm:rleb-gppa-physical} 的条件可直接验证。
取完整允许关系 $\mathcal J_v=J^v_{\lambda F}$，其中 $F=\partial f$。
前述全域唯一性给出 (G1)，而原关系的真实残差为
$r_F(t,s)=|t|=d((t,s),S)$，故 (G4) 的误差界可取 $\psi(q)=q$。
对任意输入 $x=(t,s)$、完整输出 $y=(u,s)$ 及最近零点 $p=(0,s)$，有
\[
 \delta=\|v(x)-v(p)\|=|t|^3,
 \qquad
 \langle v(y)-v(p),v(x)-v(y)\rangle
 =u^3(t^3-u^3)=\lambda u^4\ge0.
\]
因此 (G2) 可取 $a(r)=r^3$、$\omega_v(q)=q$，使用的正是同一个最近零点 $p$。
对全部完整转移，第二坐标不变，第一坐标同号且 $|u|\le|t|$，故
\[
 \|y-x\|^3=|u-t|^3
 \le|u^3-t^3|=\|v(y)-v(x)\|.
\]
于是 (G3) 可取 $h(t,q)=q^{1/3}$，相应的核增量与物理位移模为
\[
 K_v(r)=r^3,\qquad H_v(r)=r,\qquad
 \Theta_v(r)=\frac{r}{1-\kappa}.
\]
取 $0<R<\sqrt\lambda$、$\kappa=R^2/\lambda$ 及 $\bar t\ge R^3/\lambda$。
对 $0\le r\le R$，有
\[
 \frac{K_v(R)}\lambda\le\bar t,\qquad
 \psi\!\left(\frac{K_v(r)}\lambda\right)
 =\frac{r^3}\lambda\le\kappa r,
 \qquad \Theta_v(R)<\infty.
\]
这验证了输出评价窗与 (G5)，并给出一步界 $d(y,S)\le d(x,S)^3/\lambda$。
取 $U=\mathbb R^2$，则 $U_R=\{(t,s):|t|\le R\}$，
其中每个初值均满足严格预算~\eqref{eq:gppa-physical-budget}。
因此，即使核逆在解附近不满足 Lipschitz 界，该定理仍给出物理有限长度。

若改用二次核 $\phi_0(t,s)=(t^2+s^2)/2$，
同一个目标的普通 PPA 则给 $t_{k+1}=t_k/(1+\lambda)$。
因此这里的立方速率来自近端正则项的改变。
本例说明非线性核可以改变物理收敛率；
前述三角例则展示了非 Lipschitz 近端映射及沿解集的可和漂移。
